% apendices/referencia-keigo.tex
\chapter{\ruby{敬語}{けいご}\ruby{参照表}{さんしょうひょう} — Tabla de Referencia Keigo}

\section*{\ruby{動詞}{どうし}の\ruby{敬語形}{けいごけい} — Verbos en Keigo}

\begin{tcolorbox}[vocabbox, title={\ruby{特殊敬語動詞}{とくしゅけいごどうし} — Verbos Honoríficos/Humildes Especiales}]
\begin{tabularx}{\linewidth}{lllX}
\toprule
\textbf{\ruby{普通語}{ふつうご}} & \textbf{\ruby{尊敬語}{そんけいご}} & \textbf{\ruby{謙譲語}{けんじょうご}} & \textbf{Significado} \\ \midrule
いる & いらっしゃる & おる & estar \\
いく & いらっしゃる & \ruby{参}{まい}る & ir \\
くる & いらっしゃる & \ruby{参}{まい}る & venir \\
いう & おっしゃる & \ruby{申}{もう}す & decir \\
する & なさる & いたす & hacer \\
たべる & めしあがる & いただく & comer \\
のむ & めしあがる & いただく & beber \\
もらう & — & いただく & recibir \\
あげる & — & \ruby{差}{さ}し\ruby{上}{あ}げる & dar (a superior) \\
くれる & くださる & — & dar (a mí) \\
みる & ご\ruby{覧}{らん}になる & \ruby{拝見}{はいけん}する & ver \\
しる & ご\ruby{存知}{ぞんじ} & \ruby{存}{ぞん}じる & saber \\
きく & — & \ruby{拝聴}{はいちょう}する & escuchar \\
たずねる & — & \ruby{伺}{うかが}う & visitar/preguntar \\
もらう(\ruby{注文}{ちゅうもん}) & — & \ruby{承}{うけたまわ}る & recibir pedido \\
\bottomrule
\end{tabularx}
\end{tcolorbox}

\section*{\ruby{パターン}{パターン}による\ruby{敬語}{けいご} — Keigo por Patrones}

\begin{tcolorbox}[vocabbox, title={お〜になる / お〜する — Patrones Generales}]
\begin{tabularx}{\linewidth}{lXX}
\toprule
\textbf{Patrón} & \textbf{Uso} & \textbf{Ejemplos} \\ \midrule
お + V-stem + になる & \ruby{尊敬語}{そんけいご}: acciones del superior & お\ruby{帰}{かえ}りになる・お\ruby{読}{よ}みになる \\
お + V-stem + する & \ruby{謙譲語}{けんじょうご}: propias acciones & お\ruby{持}{も}ちする・お\ruby{届}{とど}けする \\
ご + N(kango) + になる & \ruby{尊敬語}{そんけいご}: sustantivos kango & ご\ruby{覧}{らん}になる・ご\ruby{使用}{しよう}になる \\
ご + N(kango) + する & \ruby{謙譲語}{けんじょうご}: sustantivos kango & ご\ruby{案内}{あんない}する・ご\ruby{説明}{せつめい}する \\
〜させていただく & Permiso humilde & \ruby{確認}{かくにん}させていただく \\
\bottomrule
\end{tabularx}
\end{tcolorbox}

\section*{\ruby{場面別}{ばめんべつ}\ruby{敬語}{けいご} — Keigo por Situación}

\begin{tcolorbox}[vocabbox, title={\ruby{電話}{でんわ}・\ruby{職場}{しょくば}の\ruby{定番表現}{ていばんひょうげん} — Expresiones Estándar de Trabajo}]
\begin{tabularx}{\linewidth}{XX}
\toprule
\textbf{Situación} & \textbf{Expresión keigo} \\ \midrule
Contestar el teléfono & はい、〜\ruby{会社}{かいしゃ}でございます。 \\
«El Sr. X no está» & X はただいま\ruby{外出}{がいしゅつ}しております。 \\
«¿Puedo tomar un mensaje?» & ご\ruby{伝言}{でんごん}を\ruby{承}{うけたまわ}りましょうか？ \\
Presentarse & \ruby{私}{わたし}は〜と\ruby{申}{もう}します。 \\
Recibir orden & \ruby{承知}{しょうち}いたしました。 \\
Disculparse (leve) & \ruby{失礼}{しつれい}いたしました。 \\
Pedir espera & \ruby{少々}{しょうしょう}お\ruby{待}{ま}ちください。 \\
Entregar documento & こちらをご\ruby{覧}{らん}ください。 \\
Recibir y leer doc. & \ruby{拝見}{はいけん}いたします。 \\
Confirmar & ご\ruby{確認}{かくにん}いただけますでしょうか？ \\
\bottomrule
\end{tabularx}
\end{tcolorbox}

\section*{\ruby{敬語}{けいご}の\ruby{注意点}{ちゅうしょく} — Puntos Clave}

\begin{tcolorbox}[grammarbox]
\begin{enumerate}
  \item \textbf{\ruby{尊敬語}{そんけいご}は\ruby{自分}{じぶん}に\ruby{使}{つか}わない}: Nunca usar 尊敬語 para describir las propias acciones.\\
  ✗「私がおっしゃいました。」 \quad ✓「私が申しました。」
  \item \textbf{\ruby{謙譲語}{けんじょうご}は\ruby{相手}{あいて}に\ruby{使}{つか}わない}: Nunca usar 謙譲語 para las acciones del superior.\\
  ✗「先生が参りました。」 \quad ✓「先生がいらっしゃいました。」
  \item \textbf{二\ruby{重敬語}{じゅうけいご}に\ruby{注意}{ちゅうい}}: Evitar la doble keigo redundante.\\
  ✗「おっしゃられました。」 \quad ✓「おっしゃいました。」
  \item \textbf{\ruby{内外}{うちそと}の\ruby{区別}{くべつ}}: Al hablar con externos, humillar al propio jefe.\\
  「\ruby{部長}{ぶちょう}の\ruby{山田}{やまだ}はただいま\ruby{不在}{ふざい}でございます。」 (humillar al jefe ante cliente)
\end{enumerate}
\end{tcolorbox}

\section*{\ruby{丁寧語}{ていねいご}のバリエーション — Variantes de Cortesía}

\begin{tcolorbox}[vocabbox, title={です・ます → でございます — Escala de cortesía}]
\begin{tabularx}{\linewidth}{lXX}
\toprule
\textbf{Nivel} & \textbf{Ejemplo} & \textbf{Contexto} \\ \midrule
\ruby{普通体}{ふつうたい} & これは\ruby{本}{ほん}だ。 & amigos, familia \\
\ruby{丁寧体}{ていねいたい} & これは\ruby{本}{ほん}です。 & conocidos, colegas \\
\ruby{丁重語}{ていちょうご} & これは\ruby{本}{ほん}でございます。 & clientes, muy formal \\
\bottomrule
\end{tabularx}
\end{tcolorbox}
